\section{The Hermitian determinant normal slice as a Lorentzian host}
\label{hnsp-section}

This independently extractable section restates one already verified
finite-dimensional identity and constructs exact physical models from it.
When it is placed after the nine-dimensional observation construction, the
notation agrees by
\(C_{A_2}=C\), \(R_P=R\), \(H_0=uu^*\),
\(Q_\perp=Q_u\), and the same map
\(\jmath_P(G)=R_PGR_P^*\).  The positive conclusion is not
merely that the identity resembles Minkowski space: the affine normal slice
has a canonically polarized Lorentz metric, the displayed \(SL_2(\mathbb C)\)
congruence subgroup induces the spin double cover of the proper orthochronous
Lorentz group, and after one explicit choice of length scale it supports a completely
specified Klein--Gordon theory and finite Gaussian regulators.  The same host
also supports classical Yang--Mills theory for any separately chosen compact
gauge group.  In addition, the standard metrics entering the Moore--Penrose
section select one exact compact \(SU(2)\)-type representation internal to the
normal data; promoting its pointwise action to a local gauge symmetry gives
the gauge--scalar model below and is an additional construction.  The proofs
identify exactly which structures come from the determinant and which are
additional physical choices.

\subsection{Fixed matrices and the exact normal metric}

Put
\begin{equation}
 u=\begin{pmatrix}1\\1\\1\end{pmatrix},\qquad
 P=\begin{pmatrix}1&-1&0\\0&1&-1\end{pmatrix},\qquad
 C_{A_2}=PP^*=\begin{pmatrix}2&-1\\-1&2\end{pmatrix},
 \label{hnsp-fixed-data}
\end{equation}
and use the standard Hermitian products in these displayed bases.  The
standard-metric Moore--Penrose right inverse is
\begin{equation}
 R_P=P^*(PP^*)^{-1}
 =\frac13\begin{pmatrix}2&1\\-1&1\\-1&-2\end{pmatrix}.
 \label{hnsp-moore-penrose}
\end{equation}
The denominators in \eqref{hnsp-moore-penrose} are essential: this is a
metric-selected complex-linear section, not a homomorphism splitting the
displayed integral lattices.
Set
\begin{equation}
 H_0=uu^*,\qquad
 \mathcal N_P=\{X\in\operatorname{Herm}_3(\mathbb C):Xu=0\},\qquad
 \jmath_P(G)=R_P G R_P^*.
 \label{hnsp-normal-space}
\end{equation}

\begin{proposition}[Exact identification of the normal slice]
\label{hnsp-normal-identification}
The following identities hold without a change of normalization:
\begin{equation}
 PR_P=I_2,\qquad u^*R_P=0,\qquad
 R_P^*R_P=C_{A_2}^{-1},\qquad
 R_P P=I_3-\frac13uu^*.
 \label{hnsp-right-inverse-identities}
\end{equation}
The real-linear map
\(
 \jmath_P:\operatorname{Herm}_2(\mathbb C)\to\mathcal N_P
\)
is an isomorphism.  Its inverse is
\begin{equation}
 \jmath_P^{-1}(X)=PXP^*.
 \label{hnsp-normal-inverse}
\end{equation}
For all \(G\in\operatorname{Herm}_2(\mathbb C)\) and
\(\varepsilon\in\mathbb R\),
\begin{equation}
 \det\bigl(H_0+\varepsilon\jmath_P(G)\bigr)
 =\varepsilon^2\det G.
 \label{hnsp-determinant-identity}
\end{equation}
\end{proposition}

\begin{proof}
The first three equalities in
\eqref{hnsp-right-inverse-identities} follow by direct multiplication.  The
last member is the orthogonal projection onto the row space of \(P\), which is
\(u^\perp\); direct multiplication gives the displayed matrix as well.  Hence
\(P\jmath_P(G)P^*=G\).  Conversely, if \(X=X^*\) and \(Xu=0\), then also
\(u^*X=0\), and therefore
\[
 R_P(PXP^*)R_P^*=(R_P P)X(R_P P)=X.
\]
This proves the isomorphism and its inverse.

For the determinant, put \(v_0=u/\sqrt3\), choose an orthonormal complement
\(U_\perp\), and let \(U=(v_0\ U_\perp)\) and
\(L=U_\perp^*R_P\).  Then
\[
 U^*H_0U=\operatorname{diag}(3,0,0),\qquad
 U^*\jmath_P(G)U=0\oplus LGL^*,
\]
while
\[
 |\det L|^2=\det(L^*L)=\det(R_P^*R_P)
 =\det(C_{A_2}^{-1})=\frac13.
\]
Taking the determinant in this orthonormal frame gives
\[
 \det(H_0+\varepsilon\jmath_P(G))
 =3\varepsilon^2|\det L|^2\det G
 =\varepsilon^2\det G.
\]
\end{proof}

Put \(X_P=(R_P\;u)\).  The same block computation also records the exact
position of this module inside the full nine-real-dimensional carrier:
\[
 H_0+\varepsilon\jmath_P(G)
 =X_P\begin{pmatrix}\varepsilon G&0\\0&1\end{pmatrix}X_P^*.
\]
Thus, in the full \((G,\psi,\lambda)\) block coordinates, this section uses
the strict sublocus \((\psi,\lambda)=(0,1)\).  Its field models live on the
observed four-vector section; they do not yet define dynamics for the retained
spinor \(\psi\) and scalar \(\lambda\) directions.

For \(G,K\in\operatorname{Herm}_2(\mathbb C)\), define
\begin{equation}
 \eta_2(G,K)=\frac12\bigl(\det(G+K)-\det G-\det K\bigr).
 \label{hnsp-polarized-determinant}
\end{equation}
Transport this form to \(\mathcal N_P\) by
\begin{equation}
 \eta_P(\jmath_P(G),\jmath_P(K))=\eta_2(G,K).
 \label{hnsp-transported-metric}
\end{equation}

\begin{theorem}[The determinant Hessian is exactly Minkowski]
\label{hnsp-lorentz-theorem}
For tangent vectors \(X=\jmath_P(G)\) and \(Y=\jmath_P(K)\),
\begin{equation}
 \eta_P(X,Y)=\frac12D^2(\det)_{H_0}[X,Y].
 \label{hnsp-hessian-metric}
\end{equation}
Writing
\begin{equation}
 G(t,x,y,z)=
 \begin{pmatrix}t+y&x+iz\\x-iz&t-y\end{pmatrix},
 \label{hnsp-pauli-coordinates}
\end{equation}
Relative to the convention
\(
 \left(\begin{smallmatrix}
 t+x_3&x_1-i x_2\\ x_1+i x_2&t-x_3
 \end{smallmatrix}\right)
\)
used in the observation construction, the exact coordinate dictionary is
\((x_1,x_2,x_3)=(x,-z,y)\).  Thus the two displayed Pauli conventions differ
only by this stated spatial relabelling and sign, not by a change of Lorentz
form.
one has
\begin{equation}
 \eta_2(G(t,x,y,z),G(t',x',y',z'))
 =tt'-xx'-yy'-zz'.
 \label{hnsp-minkowski-form}
\end{equation}
Thus \((\mathcal N_P,\eta_P)\) has signature \((1,3)\).  Its distinguished
future cone is
\begin{equation}
 \mathcal C_P^+=\jmath_P\bigl(\operatorname{Herm}_2(\mathbb C)_{\geq0}\bigr)
 =\{\jmath_P(G(t,\mathbf x)):t\geq|\mathbf x|\}.
 \label{hnsp-future-cone}
\end{equation}
Its nonzero null boundary consists exactly of the images of rank-one positive
semidefinite matrices.
\end{theorem}

\begin{proof}
Proposition~\ref{hnsp-normal-identification}, applied to \(sG+rK\), gives
\[
 \det(H_0+s\jmath_P(G)+r\jmath_P(K))=\det(sG+rK).
\]
The coefficient of \(sr\) on the right is \(2\eta_2(G,K)\).  Taking one
derivative in each variable proves \eqref{hnsp-hessian-metric}.  Direct
expansion of \eqref{hnsp-pauli-coordinates} gives
\[
 \det G(t,x,y,z)=t^2-x^2-y^2-z^2,
\]
and polarization gives \eqref{hnsp-minkowski-form}.  A two-by-two Hermitian
matrix is positive semidefinite precisely when its trace is nonnegative and
its determinant is nonnegative, which here is
\(t\geq0\) and \(t^2\geq|\mathbf x|^2\).  On the nonzero boundary its
determinant vanishes and its rank is one.
\end{proof}

The corresponding tetrad inside \(\operatorname{Herm}_3(\mathbb C)\) is
completely explicit.  If
\(E_0=\jmath_P(G(1,0,0,0))\),
\(E_1=\jmath_P(G(0,1,0,0))\),
\(E_2=\jmath_P(G(0,0,1,0))\), and
\(E_3=\jmath_P(G(0,0,0,1))\), then
\begin{align}
 E_0&=\frac19\begin{pmatrix}5&-1&-4\\-1&2&-1\\-4&-1&5\end{pmatrix},&
 E_1&=\frac19\begin{pmatrix}4&1&-5\\1&-2&1\\-5&1&4\end{pmatrix},
 \label{hnsp-explicit-tetrad-real}\\
 E_2&=\frac13\begin{pmatrix}1&-1&0\\-1&0&1\\0&1&-1\end{pmatrix},&
 E_3&=\frac13\begin{pmatrix}0&i&-i\\-i&0&i\\i&-i&0\end{pmatrix}.
 \label{hnsp-explicit-tetrad-imaginary}
\end{align}
Every row and column of each \(E_\mu\) sums to zero, and
\begin{equation}
 \eta_P(E_\mu,E_\nu)=\operatorname{diag}(1,-1,-1,-1)_{\mu\nu}.
 \label{hnsp-tetrad-metric}
\end{equation}
These identities follow by multiplying the matrices in
\eqref{hnsp-moore-penrose} and then applying
\eqref{hnsp-minkowski-form}; no orthonormalization of the marked lattice is
being suppressed.

The metric is a genuine normal Hessian.  It is not a positive energy form:
it takes both signs.  Energy appears only after choosing a time function and
passing to a Hamiltonian formulation below.

Because
\(
 \det(H_0+\jmath_P(G))=\det G
\)
on the entire affine slice, not just to second order, its restriction there is
quadratic.  Hence the polarized Hessian \(\eta_P\) is constant at every point
of the affine slice; affine translations really are isometries of this metric.

\subsection{Exact Lorentz and affine symmetry}

For \(A\in SL_2(\mathbb C)\), define
\begin{equation}
 \widehat A_P=\frac13uu^*+R_P A P\in SL_3(\mathbb C).
 \label{hnsp-lifted-spin-action}
\end{equation}

\begin{theorem}[The spin action inside the normal slice]
\label{hnsp-spin-theorem}
The matrix \(\widehat A_P\) is invertible, with
\begin{equation}
 \widehat A_P^{-1}=\frac13uu^*+R_P A^{-1}P,
 \qquad \det\widehat A_P=\det A=1.
 \label{hnsp-spin-inverse}
\end{equation}
It fixes \(u\) and satisfies
\begin{equation}
 \widehat A_P H_0\widehat A_P^*=H_0,
 \qquad
 \widehat A_P\jmath_P(G)\widehat A_P^*
 =\jmath_P(AGA^*).
 \label{hnsp-spin-congruence}
\end{equation}
Consequently \(SL_2(\mathbb C)\) acts on \(\mathcal N_P\) by future-cone
preserving \(\eta_P\)-isometries.  The kernel is \(\{I_2,-I_2\}\), and the
induced map
\begin{equation}
 PSL_2(\mathbb C)\xrightarrow{\ \sim\ }SO^+(1,3)
 \label{hnsp-spin-cover}
\end{equation}
is an isomorphism of real Lie groups.
\end{theorem}

The kernel \(\{\pm I_2\}\) is a statement about the normal-vector action.
On the full \((G,\psi,\lambda)\) carrier the same embedded
\(SL_2(\mathbb C)\)-action is faithful, because \(-I_2\) fixes \(G\) and
\(\lambda\) but sends \(\psi\mapsto-\psi\).

\begin{proof}
Let \(Q_0=uu^*/3\) and \(Q_\perp=R_P P=I_3-Q_0\).  The relations
\(Q_0R_P=0\), \(PQ_0=0\), and \(PR_P=I_2\) show that the two matrices in
\eqref{hnsp-spin-inverse} multiply to \(Q_0+Q_\perp=I_3\).  In the ordered
decomposition \(\mathbb Cu\oplus R_P\mathbb C^2\), the map is
\(1\oplus A\), proving its determinant.  The same relations prove
\eqref{hnsp-spin-congruence}.  Since
\(\det(AGA^*)=|\det A|^2\det G=\det G\), the action preserves
\(\eta_P\); congruence preserves positive semidefiniteness, hence the chosen
future cone.

If \(AGA^*=G\) for every Hermitian \(G\), taking \(G=I_2\) makes \(A\)
unitary, after which \(A\) commutes with every Hermitian matrix.  Hence
\(A\) is scalar, and determinant one gives \(A=\pm I_2\).  The differential
of the induced map has zero kernel: if
\(XG+GX^*=0\) for every Hermitian \(G\), then \(G=I_2\) makes \(X\)
skew-Hermitian and the remaining equations make it scalar; trace zero then
gives \(X=0\).  Domain and target Lie algebras both have real dimension six,
so the image contains a neighborhood of the identity.  It is therefore an
open subgroup of the connected group \(SO^+(1,3)\), hence the whole group.
\end{proof}

Translations are also intrinsic to the affine slice.  If
\(B\in\operatorname{Herm}_2(\mathbb C)\), define
\begin{equation}
 (B,A)\cdot\bigl(H_0+\jmath_P(G)\bigr)
 =H_0+\jmath_P(B+AGA^*).
 \label{hnsp-poincare-action}
\end{equation}
With multiplication
\(
 (B,A)(C,D)=(B+ACA^*,AD),
\)
this is the spin double cover of the proper orthochronous Poincar\'e action.
Its kernel is \(\{(0,I_2),(0,-I_2)\}\).
Under the dimensional chart \eqref{hnsp-dimensional-chart} introduced below,
a physical translation by \(b\in\mathbb R^{1,3}\) is the element
\(B=G(b/\ell)\).

\begin{proposition}[Ambient \(SU(3)\) is not an internal color action]
\label{hnsp-su3-distinction}
The subgroup of \(SU(3)\) preserving \(H_0\) is
\begin{equation}
 \operatorname{Stab}_{SU(3)}(H_0)
 =\{g\in SU(3):gu=e^{i\theta}u\}
 \cong S(U(1)\times U(2)).
 \label{hnsp-su3-stabilizer}
\end{equation}
It preserves \(\mathcal N_P\).  A general element of \(SU(3)\) moves
\(H_0\) and hence moves the chosen affine normal slice to a different one.
Furthermore, \(\widehat A_P\) in \eqref{hnsp-lifted-spin-action} is unitary
exactly when
\begin{equation}
 A^*C_{A_2}^{-1}A=C_{A_2}^{-1};
 \label{hnsp-compact-spin-subgroup}
\end{equation}
this compact subgroup is conjugate to \(SU(2)\), whereas generic Lorentz
boosts are not unitary.  Thus the Lorentz spin action supplied by the normal
slice is not the ambient \(SU(3)\) orbit action and does not itself define an
internal \(SU(3)\) gauge symmetry.
\end{proposition}

\begin{proof}
The equality \(gH_0g^*=H_0\) is equivalent to equality of the rank-one
lines and norms, hence to \(gu=e^{i\theta}u\).  The orthogonal complement is
then invariant, giving the displayed block stabilizer and preservation of
Hermitian operators supported on \(u^\perp\).  On
\(R_P\mathbb C^2=u^\perp\), the standard norm pulls back to the form
\(R_P^*R_P=C_{A_2}^{-1}\).  Therefore \(\widehat A_P\) is unitary precisely
when \eqref{hnsp-compact-spin-subgroup} holds.  Positive-definite congruence
identifies this group with \(SU(2)\).
\end{proof}

The proposition concerns the congruence action of the ambient \(SU(3)\) on
\(\operatorname{Herm}_3(\mathbb C)\).  It is compatible with, but distinct
from, the left \(SU(3)\)-action on
\(\operatorname{Mat}_{3\times2}(\mathbb C)\) whose invariant Gram map is
\(Z\mapsto Z^*Z\): the source spaces and actions are different, and no
identification of them is being made here.

\subsection{The compact gauge--scalar carrier selected by the standard metrics}

The Lorentz congruence group is noncompact, so it cannot preserve a positive
field-space norm.  The standard metrics already used in
\eqref{hnsp-moore-penrose} do, however, select the positive matrix and compact
subgroup
\begin{equation}
 W_P:=R_P^*R_P=C_{A_2}^{-1}
 =\frac13\begin{pmatrix}2&1\\1&2\end{pmatrix},\qquad
 K_P:=\{a\in SL_2(\mathbb C):a^* W_P a=W_P\}.
 \label{hnsp-intrinsic-compact-group}
\end{equation}
Because \(W_P=C_{A_2}^{-1}\), inversion of the defining equality gives
\[
 a^*W_Pa=W_P
 \quad\Longleftrightarrow\quad
 aC_{A_2}a^*=C_{A_2}.
\]
Accordingly, \(K_P\) is exactly the Lorentz little group of the future
timelike vector \(C_{A_2}\), equivalently the observer stabilizer selected by
the standard metrics.  Its promotion to a local internal gauge symmetry below
is mathematically valid but is an additional physical interpretation.

\begin{proposition}[Exact compact representation on the normal slice]
\label{hnsp-intrinsic-compact-representation}
Conjugation \(a\mapsto W_P^{1/2}aW_P^{-1/2}\) is an isomorphism
\(K_P\cong SU(2)\).  The real bilinear form
\begin{equation}
 E_P(G,K):=\operatorname{tr}\bigl(\jmath_P(G)\jmath_P(K)\bigr)
 =\operatorname{tr}(G W_P K W_P)
 \label{hnsp-positive-normal-metric}
\end{equation}
is a positive-definite \(K_P\)-invariant inner product on
\(\operatorname{Herm}_2(\mathbb C)\).  For \(a\in K_P\), the lift
\begin{equation}
 \widehat a_P=\frac13uu^*+R_P a P
 \label{hnsp-compact-lift}
\end{equation}
lies in \(SU(3)\), fixes \(H_0\), and obeys
\begin{equation}
 \widehat a_P\jmath_P(G)\widehat a_P^*=\jmath_P(aGa^*).
 \label{hnsp-compact-equivariance}
\end{equation}
After whitening, \(\widetilde G=W_P^{1/2}GW_P^{1/2}\), the action is
ordinary unitary congruence.  Therefore, as a real \(K_P\)-module,
\begin{equation}
 \operatorname{Herm}_2(\mathbb C)
 =\mathbb R W_P^{-1}
  \oplus\{G\in\operatorname{Herm}_2(\mathbb C):
              \operatorname{tr}(W_PG)=0\}
 \cong\boldsymbol 1\oplus\boldsymbol 3.
 \label{hnsp-singlet-adjoint-decomposition}
\end{equation}
The centre \(\{\pm I_2\}\) acts trivially.  Thus the metric-selected normal
field is a singlet plus an adjoint triplet; it is neither a fundamental weak
doublet nor an \(SU(3)\) color triplet.
\end{proposition}

\begin{proof}
The defining identity in \eqref{hnsp-intrinsic-compact-group} is equivalent to
\((W_P^{1/2}aW_P^{-1/2})^*(W_P^{1/2}aW_P^{-1/2})=I_2\); determinants are one,
so this identifies \(K_P\) with \(SU(2)\).  Cyclicity of the trace and
\(R_P^*R_P=W_P\) give the second equality in
\eqref{hnsp-positive-normal-metric}.  Moreover
\(E_P(G,G)=\operatorname{tr}(\jmath_P(G)^2)\), which is positive for nonzero
Hermitian \(\jmath_P(G)\).  Equation
\eqref{hnsp-spin-congruence} and unitarity of \(\widehat a_P\) prove
invariance.  Proposition~\ref{hnsp-su3-distinction} gives
\(\widehat a_P\in U(3)\), and the block action on
\(u^\perp\oplus\mathbb Cu\) gives
\(\det\widehat a_P=\det a=1\).  The whitened formula is direct
multiplication.  Unitary congruence fixes the scalar line and acts through the
adjoint rotation representation on traceless Hermitian matrices; the central
elements act by \(G\mapsto(\pm I_2)G(\pm I_2)^*=G\).
\end{proof}

\begin{theorem}[A fully typed positive gauge--scalar realization]
\label{hnsp-intrinsic-gauge-scalar-theorem}
Let \((M,g)\) be an oriented Riemannian \(d\)-manifold and use the trivial
principal \(K_P\)-bundle.  Take
\begin{equation}
 \mathcal A\in\Omega^1(M;\mathfrak k_P),\qquad
 G\in C^\infty(M;\operatorname{Herm}_2(\mathbb C))
 \label{hnsp-intrinsic-fields}
\end{equation}
as the connection and normal scalar.  For
\(a\in C^\infty(M,K_P)\), define
\begin{align}
 G^a&=aGa^*,&
 \mathcal A^a&=a\mathcal A a^{-1}-(da)a^{-1},
 \label{hnsp-intrinsic-gauge-action}\\
 D_{\mathcal A}G&=dG+\mathcal A G+G\mathcal A^*,&
 F_{\mathcal A}&=d\mathcal A+\mathcal A\wedge\mathcal A.
 \label{hnsp-intrinsic-covariant-objects}
\end{align}
Choose on \(\mathfrak k_P\) the positive
\(\operatorname{Ad}(K_P)\)-invariant form transported from
\(-\operatorname{Re}\operatorname{tr}(XY)\) on \(\mathfrak{su}(2)\).  For
\(g_{\rm YM},Z_G>0\) and \(m_G,\lambda_4\geq0\), the Euclidean action
\begin{equation}
 \begin{split}
 S^{\rm cpt}_{g_{\rm YM},Z_G,m_G,\lambda_4}[\mathcal A,G]
 =\int_M\bigg(&\frac1{2g_{\rm YM}^2}|F_{\mathcal A}|^2
 +\frac {Z_G}2|D_{\mathcal A}G|_{E_P}^2
 +\frac{m_G^2}{2}|G|_{E_P}^2\\
 &+\frac{\lambda_4}4(\det G)^2\bigg)\,d\operatorname{vol}_g
 \end{split}
 \label{hnsp-intrinsic-positive-action}
\end{equation}
is nonnegative and gauge invariant.  The associated normal-matrix field
\begin{equation}
 H_\varepsilon=H_0+\varepsilon\jmath_P(G)
 \label{hnsp-intrinsic-normal-field}
\end{equation}
is equivariant under the lifted action
\(H_\varepsilon^a=\widehat a_P H_\varepsilon\widehat a_P^*\).  If
\(\widehat{\mathcal A}=R_P\mathcal A P\), then
\begin{equation}
 D_{\widehat{\mathcal A}}\jmath_P(G)
 =\jmath_P(D_{\mathcal A}G),\qquad
 D_{\widehat{\mathcal A}}H_0=0.
 \label{hnsp-intrinsic-connection-intertwining}
\end{equation}
On \(M=\mathbb R^d\), the scalar Hessian at
\((\mathcal A,G)=(0,0)\) is the nonnegative self-adjoint operator
\begin{equation}
 \begin{aligned}
  \mathcal K_{Z_G,m_G}
   &=(-Z_G\Delta+m_G^2)\otimes I_{\operatorname{Herm}_2},\\
  \operatorname{Dom}(\mathcal K_{Z_G,m_G})
   &=H^2(\mathbb R^d;\operatorname{Herm}_2(\mathbb C)),\\
  \operatorname{Spec}(\mathcal K_{Z_G,m_G})
   &=[m_G^2,\infty).
 \end{aligned}
 \label{hnsp-intrinsic-scalar-hessian}
\end{equation}
Thus the fixed normal data plus the standard metrics select the compact
representation \(K_P\curvearrowright\operatorname{Herm}_2(\mathbb C)\) and
its positive fibre metric.  After the additional choice to gauge that
pointwise action, the formulas above give an exact compact gauge--scalar
model, but they still do not select \(m_G,Z_G,g_{\rm YM},\lambda_4\), a quantum
state, or an interacting gap.
\end{theorem}

\begin{proof}
Direct substitution gives
\(D_{\mathcal A^a}G^a=a(D_{\mathcal A}G)a^*\) and
\(F_{\mathcal A^a}=aF_{\mathcal A}a^{-1}\).  Invariance of the two positive
inner products and
\(\det(aGa^*)=|\det a|^2\det G=\det G\) prove gauge invariance; every term in
\eqref{hnsp-intrinsic-positive-action} is nonnegative.  Equation
\eqref{hnsp-compact-equivariance} gives equivariance of
\eqref{hnsp-intrinsic-normal-field}.  Using \(PR_P=I_2\), \(Pu=0\), and
\(u^*R_P=0\), direct multiplication proves
\eqref{hnsp-intrinsic-connection-intertwining}.  Fourier transformation gives
the self-adjoint operator, its full domain, and its spectrum in
\eqref{hnsp-intrinsic-scalar-hessian}.  Finally, \((\det G)^2\) is homogeneous
of degree four, so it has zero Hessian at \(G=0\).
\end{proof}

\begin{corollary}[The selected normal field cannot screen a fundamental centre charge]
\label{hnsp-centre-screening-corollary}
Every tensor power and every polynomial composite of the
\(\boldsymbol 1\oplus\boldsymbol 3\) normal representation has trivial
action by \(-I_2\in K_P\).  The fundamental two-dimensional representation
has action \(-I_2\).  Therefore the field \(G\) in
Theorem~\ref{hnsp-intrinsic-gauge-scalar-theorem} cannot by itself screen a
fundamental centre charge; such a model requires an additional
centre-nontrivial matter field and its independently specified kinetic,
mass, and interaction terms.
This non-screening statement is confined to the \(G\)-only normal section.
The full \((G,\psi,\lambda)\) carrier already contains the centre-nontrivial
spinor \(\psi\); what remains missing is a specified kinetic and interaction
theory for that retained field.
\end{corollary}

\begin{proof}
The centre acts trivially on \(G\), hence trivially on every tensor power and
subquotient generated from it.  It acts as minus the identity in the
fundamental representation, so that representation cannot occur in this
tensor category.
\end{proof}

\subsection{The candidate cusp matrix is an exact spacelike vector}

The preceding construction applies to every Hermitian two-by-two matrix.  For
the particular inverse-dual matrix retained in the cusp calculation, let
\begin{equation}
 G_{\eta,T}=\begin{pmatrix}2L_T&2m_T\\2m_T&-q_T/3\end{pmatrix},
 \qquad D_T=L_T q_T+6m_T^2>0.
 \label{hnsp-candidate-matrix}
\end{equation}

\begin{proposition}[Exact candidate-specific Lorentz coordinates]
\label{hnsp-candidate-spacelike}
In the coordinates \eqref{hnsp-pauli-coordinates},
\begin{equation}
 (t_T,x_T,y_T,z_T)
 =\left(L_T-\frac{q_T}{6},\ 2m_T,\ L_T+\frac{q_T}{6},\ 0\right),
 \label{hnsp-candidate-coordinates}
\end{equation}
and
\begin{equation}
 \eta_2(G_{\eta,T},G_{\eta,T})
 =\det G_{\eta,T}=-\frac23D_T<0.
 \label{hnsp-candidate-spacelike-norm}
\end{equation}
If \(L_T/T\to1\), \(q_T/T\to1\), and \(m_T/T\to0\), then
\begin{equation}
 \frac1T(t_T,x_T,y_T,z_T)
 \longrightarrow\left(\frac56,0,\frac76,0\right).
 \label{hnsp-candidate-direction-limit}
\end{equation}
Thus this particular candidate vector is spacelike for every \(T\).  It is
not a positive-mass four-momentum.  This statement concerns the indefinite
inverse-dual matrix \(G_{\eta,T}\); it does not concern the different positive
Gram momentum \(p_\alpha(Z)=\alpha Z^*Z\) used in the retained-\(c\) worldline
construction.  The latter is positive semidefinite by definition, and its
image is all of \(\operatorname{Herm}_2(\mathbb C)_{\geq0}\): for every
\(G\succeq0\), the block matrix
\(
 Z=\alpha^{-1/2}
 \left(\begin{smallmatrix}G^{1/2}\\0_{1\times2}\end{smallmatrix}\right)
 \in\operatorname{Mat}_{3\times2}(\mathbb C)
\)
satisfies \(p_\alpha(Z)=G\).  Hence
\[
 G_{\eta,T}\notin\operatorname{im}p_\alpha
 \qquad(T>0,\ \alpha>0).
\]
Positive rescaling preserves the determinant sign and \(GL_2(\mathbb C)\)
congruence preserves inertia, so neither operation repairs this mismatch.
The Gram momentum lies on \(\det p_\alpha=m_0^2c^2\) only after that
mass-shell constraint and the parameters \(\alpha,c,m_0\) are supplied.
Interpreting \(G_{\eta,T}\) as a field value, a normal deformation, or a
spacelike wave-vector requires an additional map; the determinant calculation
alone does not choose among those interpretations.
\end{proposition}

\begin{proof}
Solving \(t+y=2L_T\), \(t-y=-q_T/3\), and
\(x+iz=2m_T\) gives \eqref{hnsp-candidate-coordinates}.  Substitution gives
\[
 t_T^2-x_T^2-y_T^2-z_T^2
 =(2L_T)(-q_T/3)-4m_T^2=-\frac23D_T.
\]
The limit is componentwise.
\end{proof}

\begin{theorem}[Exact determinant-level orbit geometry]
\label{hnsp-determinant-orbit-geometry}
Fix \(\varepsilon>0\) and write
\[
 \mathbb M_\varepsilon
 :=H_0+\varepsilon\jmath_P\bigl(\operatorname{Herm}_2(\mathbb C)\bigr).
\]
For every \(r>0\), the future timelike determinant shell
\begin{equation}
 \mathscr P^+_{r,\varepsilon}
 :=\{H_0+\varepsilon\jmath_P(G):G\succ0,\ \det G=r^2\}
 \cong SL_2(\mathbb C)/SU(2)
 \label{hnsp-future-orbit}
\end{equation}
is one orbit of the lifted congruence action.  It equals
\[
 \{H\in\mathbb M_\varepsilon:H\succ0,\
   \det H=\varepsilon^2r^2\}.
\]
The nonzero future-null locus
\begin{equation}
 \mathscr N^+_\varepsilon
 :=\{H_0+\varepsilon\jmath_P(G):
       G\succeq0,\ \operatorname{rank}G=1\}
 \label{hnsp-null-orbit}
\end{equation}
is also one \(SL_2(\mathbb C)\)-orbit; each of its matrices is positive
semidefinite of ambient rank two.  Finally, with
\(r_T=\sqrt{2D_T/3}\), the candidate-specific spacelike shell is
\begin{equation}
 \mathscr S_{T,\varepsilon}
 :=\{H_0+\varepsilon\jmath_P(G):\det G=-r_T^2\}
 \cong SL_2(\mathbb C)/SU(1,1),
 \label{hnsp-spacelike-orbit}
\end{equation}
and every matrix in it has inertia \((2,1)\).  Consequently the normal slice
contains every timelike mass shell and the future null cone after a unit
dictionary is supplied, but it selects neither \(r\) nor a physical mass.
\end{theorem}

\begin{proof}
If \(G\succ0\) and \(\det G=r^2\), then
\(A=(G/r)^{1/2}\) lies in \(SL_2(\mathbb C)\) and
\(A(rI_2)A^*=G\).  The stabilizer of \(rI_2\) is \(SU(2)\), proving
\eqref{hnsp-future-orbit}; positivity and the ambient determinant follow from
the block decomposition in the proof of
Proposition~\ref{hnsp-normal-identification}.  Every nonzero positive
rank-one \(G\) is \(vv^*\).  Choose \(w\in\mathbb C^2\) with
\(\det(v\ w)=1\); then \(A=(v\ w)\in SL_2(\mathbb C)\) and
\(A\operatorname{diag}(1,0)A^*=vv^*\).  This proves null transitivity, and the
orthogonal block decomposition gives ambient rank \(1+1=2\).

Let \(J=\operatorname{diag}(1,-1)\).  A Hermitian \(G\) with
\(\det G=-r_T^2\) has inertia \((1,1)\), so Sylvester congruence gives
\(G=B(r_T J)B^*\) for some \(B\in GL_2(\mathbb C)\).  Determinants give
\(|\det B|=1\).  If \(\det B=e^{i\theta}\), then
\(D=\operatorname{diag}(e^{-i\theta},1)\) stabilizes \(J\), and
\(A=BD\) has determinant one.  Thus the action is transitive, and the
stabilizer is
\[
 \{A\in SL_2(\mathbb C):AJA^*=J\}=SU(1,1).
\]
Equation~\eqref{hnsp-candidate-spacelike-norm} puts the retained candidate on
this shell.  The block decomposition gives one positive eigenvalue from
\(H_0\) and one of each sign from \(G\), proving the inertia claim.
\end{proof}

\subsection{Units and the minimal continuum field model}

All entries in \eqref{hnsp-fixed-data}--\eqref{hnsp-normal-space} are
dimensionless.  A physical spacetime requires a scale.  For \(\ell>0\), set
\begin{equation}
 \kappa_\ell(x^0,x^1,x^2,x^3)
 =H_0+\jmath_P\!\left(
 G\!\left(\frac{x^0}{\ell},\frac{x^1}{\ell},
           \frac{x^2}{\ell},\frac{x^3}{\ell}\right)\right).
 \label{hnsp-dimensional-chart}
\end{equation}
Equip the affine slice with the constant metric \(g_{P,\ell}=\ell^2\eta_P\).
Then
\begin{equation}
 \kappa_\ell^*g_{P,\ell}
 =(dx^0)^2-(dx^1)^2-(dx^2)^2-(dx^3)^2,
 \qquad
 \det\kappa_\ell(x)=\ell^{-2}\eta_{\mu\nu}x^\mu x^\nu.
 \label{hnsp-dimensional-isometry}
\end{equation}
For the scaled normal displacement
\(H_0+\varepsilon\jmath_P(G(x/\ell))\), the determinant is multiplied by
\(\varepsilon^2\).  For \(\varepsilon\ne0\), the corresponding isometric
metric rescaling is \((\ell/|\varepsilon|)^2\eta_P\).  Hence
\(\varepsilon\) fixes no absolute energy or length by itself.

The exact unit dictionary is:
\begin{center}
\begin{tabular}{@{}>{\raggedright\arraybackslash}p{0.18\linewidth}
                    >{\raggedright\arraybackslash}p{0.34\linewidth}
                    >{\raggedright\arraybackslash}p{0.38\linewidth}@{}}
\toprule
symbol & role & status from the determinant data\\
\midrule
\(P,u,R_P,\jmath_P\) & marked linear algebra & fixed and dimensionless\\
\(\varepsilon\) & normal-displacement amplitude & fixed only when supplied\\
\(\ell\) & conversion from matrix coordinates to length & independent choice\\
\(x^0=ct\) & time coordinate & requires the independent constant \(c\)\\
\(m_{\mathrm{KG}}\) & inverse Compton length in the scalar model & independent choice\\
\(m_{\rm phys}=\hbar m_{\mathrm{KG}}/c\) & physical scalar mass & requires \(\hbar,c\)\\
\(L=Na\) & infrared box size and ultraviolet mesh & regulator choices\\
\(K,g_{\rm YM}\) & compact gauge group and coupling & \(K\) arbitrary; \(K_P\cong SU(2)\) is the metric-selected little group and its gauging is optional\\
\bottomrule
\end{tabular}
\end{center}

The retained worldline mass \(m_0\), the Klein--Gordon inverse length
\(m_{\mathrm{KG}}\), and the dimensionless determinant-shell radius \(r\) are
independent inputs.  Only after declaring that the worldline and field theory
describe the same physical species may one impose the unit conversion
\[
 m_{\mathrm{KG}}=\frac{m_0c}{\hbar},
 \qquad
 \Delta E_{\mathrm{free}}=\hbar c\,m_{\mathrm{KG}}=m_0c^2.
\]
No determinant identity derives this identification.

Use natural units \(c=\hbar=1\) in the remaining formulas.  The minimal
field placed on the Lorentzian host is a real scalar
\(\phi\in C_c^\infty(\mathbb R^{1,3};\mathbb R)\) with action
\begin{equation}
 S_{m_{\mathrm{KG}}}[\phi]=\frac12\int_{\mathbb R^{1,3}}
 \left(\partial_\mu\phi\,\partial^\mu\phi-m_{\mathrm{KG}}^2\phi^2\right)d^4x,
 \qquad m_{\mathrm{KG}}\geq0.
 \label{hnsp-klein-gordon-action}
\end{equation}
Its Euler--Lagrange operator is \(\Box+m_{\mathrm{KG}}^2\).  The Poincar\'e action
\eqref{hnsp-poincare-action} acts on fields by pullback.  The action, the
stress-energy tensor
\begin{equation}
 T_{\mu\nu}=\partial_\mu\phi\partial_\nu\phi
 -\frac12\eta_{\mu\nu}
  (\partial_\alpha\phi\partial^\alpha\phi-m_{\mathrm{KG}}^2\phi^2),
 \label{hnsp-stress-energy}
\end{equation}
and the mass shell \(p_\mu p^\mu=m_{\mathrm{KG}}^2\) are invariant.  Smeared fields
\(\phi(f)=\int\phi f\,d^4x\) are the basic linear observables.

On \(x^0=0\), smooth compactly supported initial data
\((\phi_0,\pi_0)\) carry the symplectic form
\begin{equation}
 \sigma((\phi,\pi),(\psi,\rho))
 =\int_{\mathbb R^3}(\phi\rho-\pi\psi)\,d^3x
 \label{hnsp-scalar-symplectic-form}
\end{equation}
and energy
\begin{equation}
 E_{m_{\mathrm{KG}}}(\phi,\pi)=\frac12\int_{\mathbb R^3}
 (|\pi|^2+|\nabla\phi|^2+m_{\mathrm{KG}}^2|\phi|^2)\,d^3x.
 \label{hnsp-scalar-energy}
\end{equation}
For \(m_{\mathrm{KG}}>0\), the energy completion is
\(H^1(\mathbb R^3)\oplus L^2(\mathbb R^3)\); for \(m_{\mathrm{KG}}=0\), it is
\(\dot H^1(\mathbb R^3)\oplus L^2(\mathbb R^3)\).  The symplectic pairing
\eqref{hnsp-scalar-symplectic-form} is exact on the dense smooth,
compactly supported data.  In the massless homogeneous completion it need not
extend continuously to every pair without an additional dual-space choice.

\begin{theorem}[Exact free-field spectral alternative]
\label{hnsp-free-field-gap}
Let
\begin{equation}
 K_{m_{\mathrm{KG}}}=-\Delta+m_{\mathrm{KG}}^2,
 \quad \operatorname{Dom}(K_{m_{\mathrm{KG}}})=H^2(\mathbb R^3),
 \quad \omega_{m_{\mathrm{KG}}}=K_{m_{\mathrm{KG}}}^{1/2}
 \label{hnsp-one-particle-operator}
\end{equation}
on \(L^2(\mathbb R^3)\).  On the bosonic Fock space
\(\mathscr F_s(L^2(\mathbb R^3))\), define the vacuum-subtracted Hamiltonian
\begin{equation}
 \mathbb H_{m_{\mathrm{KG}}}=d\Gamma(\omega_{m_{\mathrm{KG}}}).
 \label{hnsp-free-fock-hamiltonian}
\end{equation}
Explicitly, on the orthogonal decomposition
\(\mathscr F_s=\bigoplus_{n=0}^\infty L^2(\mathbb R^3)^{\otimes_s n}\), its
domain consists of the vectors \(\Psi=(\Psi_n)\) for which each \(\Psi_n\)
belongs to the domain of
\(
 \sum_{j=1}^n I\otimes\cdots\otimes\omega_{m_{\mathrm{KG}}}\otimes\cdots\otimes I
\)
and the squared norms of those images are summable.
For every \(m_{\mathrm{KG}}\geq0\), the Fock vacuum is the unique zero-energy vector.  Its
excitation gap is exactly
\begin{equation}
 \inf\operatorname{Spec}(\mathbb H_{m_{\mathrm{KG}}}|_{\Omega^\perp})=m_{\mathrm{KG}}.
 \label{hnsp-free-gap-value}
\end{equation}
Thus the same determinant normal geometry admits a positive gap of any chosen
value \(m_{\mathrm{KG}}>0\) and a gapless model at \(m_{\mathrm{KG}}=0\).
\end{theorem}

\begin{proof}
Fourier transform identifies \(K_{m_{\mathrm{KG}}}\) with multiplication by
\(|p|^2+m_{\mathrm{KG}}^2\) and \(\omega_{m_{\mathrm{KG}}}\) with multiplication by
\(\sqrt{|p|^2+m_{\mathrm{KG}}^2}\).  The latter has zero kernel even when
\(m_{\mathrm{KG}}=0\), since
the singleton \(p=0\) has Lebesgue measure zero.  Hence the only Fock vector
annihilated by \(d\Gamma(\omega_{m_{\mathrm{KG}}})\) is the vacuum.  On the
\(n\)-particle sector, the Hamiltonian is the sum of \(n\) copies of
\(\omega_{m_{\mathrm{KG}}}\), so the nonvacuum spectral bottom is at least
\(m_{\mathrm{KG}}\).  The one-particle sector is reducing and has spectral
bottom \(m_{\mathrm{KG}}\), proving equality.  For \(m_{\mathrm{KG}}=0\),
normalized Fourier-space functions supported in balls \(|p|<1/n\) have
energy expectation at most \(1/n\), which makes the absence of a positive gap
explicit.
\end{proof}

There is no contradiction between the absence of a positive kinematic
threshold on the unrestricted Gram carrier and the free Fock gap
\(m_{\mathrm{KG}}\): the former allows the rescaling \(Z\mapsto rZ\), whereas
the latter first fixes an external mass coefficient and quantizes the chosen
free Hamiltonian.

\subsection{A complete finite regulator and its two limits}

Let \(L>0\), \(N\geq2\), \(a=L/N\), and
\(V_N=(\mathbb Z/N\mathbb Z)^3\).  Give
\(Q_{L,N}=\mathbb R^{V_N}\) the inner product
\begin{equation}
 \langle q,r\rangle_a=a^3\sum_{v\in V_N}q(v)r(v).
 \label{hnsp-lattice-inner-product}
\end{equation}
The positive periodic lattice Laplacian is
\begin{equation}
 (-\Delta_a q)(v)=a^{-2}\sum_{j=1}^3
 \bigl(2q(v)-q(v+e_j)-q(v-e_j)\bigr).
 \label{hnsp-lattice-laplacian}
\end{equation}
Let \(d\lambda_a\) be Lebesgue measure associated with the Euclidean space
\((Q_{L,N},\langle\cdot,\cdot\rangle_a)\), and let
\(\Delta_{Q_{L,N}}\) denote its coordinate Laplacian, so that
\(-\Delta_{Q_{L,N}}\) is nonnegative.  For
\(m_{\mathrm{KG}}>0\), define on \(L^2(Q_{L,N},d\lambda_a)\)
\begin{equation}
 \mathbb H_{m_{\mathrm{KG}},L,N}
 =\frac12\left(-\Delta_{Q_{L,N}}
 +\langle q,(-\Delta_a+m_{\mathrm{KG}}^2)q\rangle_a\right)
 -E_{0,m_{\mathrm{KG}},L,N},
 \label{hnsp-lattice-hamiltonian}
\end{equation}
where \(E_{0,m_{\mathrm{KG}},L,N}\) is its finite oscillator ground energy.
More precisely, before subtracting \(E_{0,m_{\mathrm{KG}},L,N}\), this is the self-adjoint
operator associated with the closed form
\begin{equation}
 \mathfrak h_{m_{\mathrm{KG}},L,N}[\Psi]
 =\frac12\int_{Q_{L,N}}
 \left(|\nabla\Psi(q)|_a^2
 +\langle q,(-\Delta_a+m_{\mathrm{KG}}^2)q\rangle_a
             |\Psi(q)|^2\right)d\lambda_a(q),
 \label{hnsp-lattice-closed-form}
\end{equation}
whose domain is
\begin{equation}
 H^1(Q_{L,N})\cap
 L^2\!\left(Q_{L,N},
  \langle q,(-\Delta_a+m_{\mathrm{KG}}^2)q\rangle_a\,
  d\lambda_a(q)\right).
 \label{hnsp-lattice-form-domain}
\end{equation}

\begin{theorem}[Exact regulated spectrum and continuum limits]
\label{hnsp-regulator-theorem}
For \(n=(n_1,n_2,n_3)\in V_N\), the squared normal-mode frequencies are
\begin{equation}
 \omega_{m_{\mathrm{KG}},L,N}(n)^2
 =m_{\mathrm{KG}}^2+\frac4{a^2}\sum_{j=1}^3
 \sin^2\!\left(\frac{\pi n_j}{N}\right).
 \label{hnsp-lattice-dispersion}
\end{equation}
The operator \(\mathbb H_{m_{\mathrm{KG}},L,N}\) is self-adjoint on the tensor-product
oscillator domain, has a unique vacuum, and has exact gap \(m_{\mathrm{KG}}\).  The
subtracted constant is
\begin{equation}
 E_{0,m_{\mathrm{KG}},L,N}
 =\frac12\sum_{n\in V_N}\omega_{m_{\mathrm{KG}},L,N}(n).
 \label{hnsp-regulated-vacuum-energy}
\end{equation}

For \(m_{\mathrm{KG}}=0\), the constant mode is a free particle and the unreduced operator
has no normalized ground state.  On the mean-zero configuration space
\begin{equation}
 Q^0_{L,N}=\{q\in Q_{L,N}:\textstyle\sum_v q(v)=0\},
 \label{hnsp-mean-zero-space}
\end{equation}
the vacuum-subtracted oscillator has a unique ground state and gap
\begin{equation}
 \Delta^0_{L,N}=\frac2a\sin\!\left(\frac\pi N\right).
 \label{hnsp-massless-regulated-gap}
\end{equation}
At fixed \(L\), for every fixed \(r\in\mathbb Z^3\),
\begin{equation}
 \omega_{m_{\mathrm{KG}},L,N}(r)^2\longrightarrow
 m_{\mathrm{KG}}^2+\frac{4\pi^2|r|^2}{L^2}\qquad(N\to\infty),
 \label{hnsp-fixed-box-limit}
\end{equation}
and
\begin{equation}
 \Delta^0_{L,N}\longrightarrow\frac{2\pi}{L}.
 \label{hnsp-fixed-box-gap-limit}
\end{equation}
After the thermodynamic limit \(L\to\infty\), the massive gap remains
\(m_{\mathrm{KG}}\), whereas the massless positive-mode gap tends to zero.
\end{theorem}

\begin{proof}
The characters
\(
 v\mapsto L^{-3/2}e^{2\pi i n\cdot v/N}
\)
diagonalize \(-\Delta_a\).  Substitution gives
\eqref{hnsp-lattice-dispersion}.  Orthogonal diagonalization of the real
configuration space turns \eqref{hnsp-lattice-hamiltonian} into a tensor sum
of one-dimensional harmonic oscillators with frequencies
\(\omega_{m_{\mathrm{KG}},L,N}(n)\).  For \(m_{\mathrm{KG}}>0\), their
minimum is attained at \(n=0\) and equals \(m_{\mathrm{KG}}\), proving the
ground-state and gap claims.

When \(m_{\mathrm{KG}}=0\), the constant coordinate has zero restoring force,
hence gives
the free-particle operator and no \(L^2\) ground state.  Removing that
coordinate leaves the smallest character with one component congruent to
\(\pm1\) and the others zero, giving
\eqref{hnsp-massless-regulated-gap}.  With \(a=L/N\), the elementary limit
\(N\sin(\pi r_j/N)\to\pi r_j\) proves
\eqref{hnsp-fixed-box-limit} and \eqref{hnsp-fixed-box-gap-limit}.  The last
claims follow by sending \(L\) to infinity.
\end{proof}

This regulator supplies a literal configuration space, state Hilbert space,
quadratic Hamiltonian, observables generated by the coordinates and momenta,
and exact fixed-mode ultraviolet and lowest-frequency infrared limits.  The
displayed limits do not assert strong- or norm-resolvent convergence,
convergence of Gaussian measures, or convergence of Fock representations.  The mass
parameter is not hidden in \(P\), \(R_P\), \(\varepsilon\), or the
determinant: it is the independent coefficient in
\eqref{hnsp-klein-gordon-action}.

\subsection{An exact classical gauge-theory completion}

Let \(K\) be any compact matrix Lie group with Lie algebra \(\mathfrak k\)
and an \(\operatorname{Ad}K\)-invariant positive inner product.  On the
trivial bundle over the Lorentzian host, take a smooth compactly supported
gauge potential \(A\in\Omega_c^1(\mathbb R^{1,3};\mathfrak k)\), with
curvature
\begin{equation}
 F_A=dA+A\wedge A.
 \label{hnsp-gauge-curvature}
\end{equation}
The gauge group of smooth maps \(\mathbb R^{1,3}\to K\) that equal the
identity off a compact set acts by
\begin{equation}
 A^g=gAg^{-1}-(dg)g^{-1},\qquad F_{A^g}=gF_A g^{-1}.
 \label{hnsp-gauge-action}
\end{equation}
For an independently chosen coupling \(g_{\mathrm{YM}}>0\), the classical
action is
\begin{equation}
 S_{K,g_{\mathrm{YM}}}[A]
 =-\frac1{4g_{\mathrm{YM}}^2}
 \int\langle F_{\mu\nu},F^{\mu\nu}\rangle\,d^4x.
 \label{hnsp-yang-mills-action}
\end{equation}
It is invariant under \eqref{hnsp-gauge-action} and the Poincar\'e action.
For a closed curve \(\gamma\) and a finite-dimensional representation
\(\rho\) of \(K\),
\begin{equation}
 W_{\gamma,\rho}(A)
 =\operatorname{tr}\rho(\operatorname{Hol}_A(\gamma))
 \label{hnsp-wilson-observable}
\end{equation}
is gauge invariant.

On the spatial regulator \(\Sigma_L=(\mathbb R/L\mathbb Z)^3\), the smooth
Hamiltonian phase space before reduction consists of
\((A,E)\in\Omega^1(\Sigma_L;\mathfrak k)^2\), with gauge action
\begin{equation}
 (A,E)^g=\bigl(gAg^{-1}-(dg)g^{-1},\,gEg^{-1}\bigr).
 \label{hnsp-hamiltonian-gauge-action}
\end{equation}
Its Gauss constraint and
energy are
\begin{equation}
 d_A^*E=0,\qquad
 \mathcal E_{K,g}(A,E)=\frac1{2g_{\mathrm{YM}}^2}
 \int_{\Sigma_L}(|E|^2+|F_A|^2)\,d^3x,
 \label{hnsp-gauge-phase-space}
\end{equation}
and the reduced classical state space is the constraint locus modulo
\(C^\infty(\Sigma_L,K)\).  At this smooth level it is a set-theoretic orbit
space, generally stratified by stabilizer type; making it a smooth symplectic
space requires Sobolev completions and hypotheses on the action.  Equations
\eqref{hnsp-gauge-action} and \eqref{hnsp-gauge-phase-space} show directly
that the constraint locus, energy, and Wilson observables descend to this
quotient: in particular,
\(d_{A^g}^*E^g=g(d_A^*E)g^{-1}\), because the Hodge star acts only on the
base and the adjoint action is orthogonal.  In temporal gauge the canonical momentum is
\(\Pi_E=g_{\mathrm{YM}}^{-2}E\).  Equivalently, the symplectic form on
tangent vectors is
\begin{equation}
 \Omega_{(A,E)}\bigl((a,e),(a',e')\bigr)
 =\frac1{g_{\mathrm{YM}}^2}\int_{\Sigma_L}
 \bigl(\langle e,a'\rangle-\langle e',a\rangle\bigr)\,d\operatorname{vol}.
 \label{hnsp-gauge-symplectic-form}
\end{equation}

\begin{proposition}[The exact nonzero-mode Gaussian gauge gap]
\label{hnsp-gauge-gaussian-gap}
At \((A,E)=(0,0)\), impose the Coulomb conditions
\(d^*A=d^*E=0\) and remove harmonic one-forms.  Relative to the unweighted
\(L^2\) configuration metric, the energy Hessian is
\(g_{\mathrm{YM}}^{-2}\Delta_1\), where the generalized squared-frequency
operator is
\begin{equation}
 \Delta_1=d^*d,
 \quad
 \operatorname{Dom}(\Delta_1)
 =H^2\Omega^1(\Sigma_L;\mathfrak k)
  \cap\ker d^*\cap(\ker\Delta_1)^\perp.
 \label{hnsp-gauge-hessian-domain}
\end{equation}
on the closed Hilbert space
\begin{equation}
 \mathscr K_{L}^{\mathrm{tr}}
 =\bigl(L^2\Omega^1(\Sigma_L;\mathfrak k)\cap\ker d^*\bigr)
  \cap(\ker\Delta_1)^\perp.
 \label{hnsp-transverse-gauge-hilbert-space}
\end{equation}
Its spectrum is
\begin{equation}
 \left\{\frac{4\pi^2|n|^2}{L^2}:n\in\mathbb Z^3\setminus\{0\}\right\},
 \label{hnsp-gauge-hessian-spectrum}
\end{equation}
with transverse-polarization and \(\dim\mathfrak k\) multiplicities.  On the
unreduced transverse Gaussian Fock space after this explicit harmonic-mode
truncation, the gap is \(2\pi/L\), which tends to zero as \(L\to\infty\).
Constant gauge transformations remain after Coulomb gauge, so this statement
does not assert that its one-particle level survives restriction to a
gauge-invariant Fock subspace.
\end{proposition}

\begin{proof}
The curvature linearization at zero is \(A\mapsto dA\).  On the Coulomb
slice, the configuration part of the energy has quadratic form
\(g_{\mathrm{YM}}^{-2}\|dA\|^2/2\), whose Hessian relative to unweighted
\(L^2\) is \(g_{\mathrm{YM}}^{-2}d^*d\) with the displayed elliptic domain.
Dividing by the kinetic metric \(g_{\mathrm{YM}}^{-2}I\) gives the
squared-frequency operator \(\Delta_1=d^*d\).  Fourier characters on the
three-torus give \eqref{hnsp-gauge-hessian-spectrum}; for each nonzero
wave-vector, the condition \(d^*A=0\) leaves its two orthogonal polarizations.
Gaussian quantization replaces each Hessian eigenvalue \(\lambda\) by the
oscillator frequency \(\sqrt\lambda\): in canonical variables a single mode
has Hamiltonian
\(
 \tfrac12(g_{\mathrm{YM}}^2P_Q^2+
 \lambda g_{\mathrm{YM}}^{-2}Q^2)
\), whose frequency is \(\sqrt\lambda\).  The smallest nonzero wave-vector
has length \(2\pi/L\), proving the result.
\end{proof}

Thus the determinant slice is an exact four-dimensional Lorentzian host for
classical gauge theory, not merely an analogy.  It does not select \(K\): the
same construction works for \(U(1)\), \(SU(2)\), \(SU(3)\), or any other
compact matrix group.  In particular, Proposition~\ref{hnsp-su3-distinction}
prevents the ambient \(SU(3)\) orbit label from being silently reused as a
local color gauge group.  Nor does the free Hessian establish the nonlinear
four-dimensional Yang--Mills mass gap; that problem concerns the physical
gauge-invariant quantum Hamiltonian after a continuum construction
\cite{JaffeWitten2000}.

A cubical mesh gives the exact holonomy regulator of Wilson
\cite{Wilson_1974}.  To type the state dependence, let \(V\) be the vertices,
choose one orientation \(E^+\) of the edges, and put
\begin{equation}
 \mathcal C_E=K^{E^+},\qquad \mathcal G_V=K^V,
 \qquad (h\cdot U)_e=h_{s(e)}U_e h_{t(e)}^{-1}.
 \label{hnsp-link-configuration}
\end{equation}
Reverse orientation means \(U_{\bar e}=U_e^{-1}\).  For a closed edge path
\(\gamma=e_1\cdots e_r\) and a nontrivial irreducible representation \(\rho\),
use the normalized Wilson observable
\begin{equation}
 w_{\gamma,\rho}(U)
 =\frac1{\dim\rho}\operatorname{tr}
   \rho(U_{e_1}\cdots U_{e_r}).
 \label{hnsp-normalized-link-wilson}
\end{equation}

\begin{proposition}[Link kinematics selects no confinement state]
\label{hnsp-link-state-countermodels}
There are two \(\mathcal G_V\)-invariant probability measures on the same
configuration space \(\mathcal C_E\) such that
\begin{equation}
 \mathbb E_{\mathrm{pg}}[w_{\gamma,\rho}]=1
 \quad\text{for every closed }\gamma,
 \qquad
 \mathbb E_{\mathrm{Haar}}[w_{\gamma,\rho}]=0
 \label{hnsp-link-state-values}
\end{equation}
whenever \(\gamma\) contains an unoriented edge exactly once in the second
formula.  Their convex mixtures give every expectation between zero and one.
Thus the link configuration and gauge action do not determine an area law or
string tension.
\end{proposition}

\begin{proof}
For the first measure, give \(K^V\) product Haar measure and push it forward
under \(h\mapsto U(h)\), where
\(U(h)_e=h_{s(e)}h_{t(e)}^{-1}\).  Left invariance of vertex Haar measure
makes the push-forward gauge invariant, and every closed holonomy telescopes
to the identity, proving the first value.  For the second, give the independent
variables \(U_e\), \(e\in E^+\), product Haar measure.  Bi-invariance makes
this measure gauge invariant.  Integrating first over an edge that occurs
exactly once inserts
\(\int_K\rho(k)\,dk=0\), because a nontrivial irreducible representation has
no invariant vector.  This proves the second value.  Convexity proves the
remaining assertion.
\end{proof}

The link system is a valid ultraviolet regulator only after a time direction,
Euclidean rotation, coupling normalization, action, and probability measure
have been chosen.  Proposition~\ref{hnsp-link-state-countermodels} shows why
the normal-slice geometry and link space alone do not supply confinement.

\subsection{Exact obstructions and the sufficient spectral bridge}

\begin{proposition}[No Lorentz-invariant positive quadratic energy]
\label{hnsp-no-positive-lorentz-energy}
There is no positive-definite quadratic form on \(\mathcal N_P\) invariant
under the full \(PSL_2(\mathbb C)\) action of
Theorem~\ref{hnsp-spin-theorem}.  Consequently the determinant Hessian cannot
itself be a nonnegative Hamiltonian.
\end{proposition}

\begin{proof}
Suppose a positive-definite norm were invariant.  In the coordinates
\eqref{hnsp-pauli-coordinates}, Lorentz boosts send
\((1,0,0,0)\) to \((\cosh s,\sinh s,0,0)\).  Invariance would keep the
positive-definite norm of this orbit constant.  Every positive-definite norm
on a finite-dimensional space is equivalent to the Euclidean norm, whereas
the Euclidean norms of these vectors diverge with \(|s|\).  This is a
contradiction.
\end{proof}

The next two models make the role of entanglement exact.  They use bipartite
tensor-product Hilbert kinematics but have opposite gap behavior.

\begin{proposition}[Entanglement determines neither positivity nor magnitude of a gap]
\label{hnsp-entanglement-countermodels}
Let
\(\Omega=(e_0\otimes e_0+e_1\otimes e_1)/\sqrt2\) in
\(\mathbb C^2\otimes\mathbb C^2\).  Then
\begin{equation}
 H_{\mathrm{gap}}=I-|\Omega\rangle\langle\Omega|
 \label{hnsp-entangled-gapped}
\end{equation}
is bounded, nonnegative, and self-adjoint, has the entangled vector
\(\Omega\) as its unique vacuum, and has gap one.

On \(\ell^2(\mathbb N_0)\otimes\ell^2(\mathbb N_0)\), put
\begin{equation}
 \Omega_0=\frac{e_0\otimes e_0+e_1\otimes e_1}{\sqrt2},\qquad
 \psi_n=\frac{e_{2n}\otimes e_{2n}+e_{2n+1}\otimes e_{2n+1}}{\sqrt2}
 \quad(n\geq1).
 \label{hnsp-entangled-basis}
\end{equation}
Define a bounded self-adjoint operator by
\begin{equation}
 H_{\mathrm{ent}}\Omega_0=0,\qquad
 H_{\mathrm{ent}}\psi_n=n^{-1}\psi_n,
 \qquad H_{\mathrm{ent}}=I\ \text{on }
 \overline{\operatorname{span}}\{\Omega_0,\psi_n:n\geq1\}^\perp.
 \label{hnsp-entangled-gapless}
\end{equation}
Put
\[
 \mathcal K_{\mathrm{ent}}
 =\overline{\operatorname{span}}\{\psi_n:n\geq1\}.
\]
Every nonzero vector in \(\mathcal K_{\mathrm{ent}}\) is entangled,
\(\mathcal K_{\mathrm{ent}}\) reduces \(H_{\mathrm{ent}}\), and
\[
 \inf\operatorname{Spec}
  (H_{\mathrm{ent}}|_{\mathcal K_{\mathrm{ent}}})=0.
\]
Thus even a closed, completely entangled reducing sector need not have a
positive energy floor.
\end{proposition}

\begin{proof}
Equation \eqref{hnsp-entangled-gapped} is an orthogonal projection with
kernel \(\mathbb C\Omega\).  The vectors in
\eqref{hnsp-entangled-basis} are mutually orthonormal and each has Schmidt
rank two.  Equation \eqref{hnsp-entangled-gapless} is a bounded real diagonal
operator in the indicated orthogonal decomposition.  Its kernel is
\(\mathbb C\Omega_0\), while \(1/n\to0\).  Under the standard
coefficient-operator identification of a bipartite tensor, a nonzero vector
\(\sum_{n\geq1}c_n\psi_n\) has, for every \(n\) with \(c_n\ne0\), two equal
nonzero diagonal entries on the coordinates \(2n,2n+1\).  Its Schmidt rank is
therefore at least two.  The displayed diagonal action also proves reduction
of \(\mathcal K_{\mathrm{ent}}\) and spectral bottom zero.
\end{proof}

Finally, let \(A_{\mathrm{phys}}\) be a lower-semibounded self-adjoint
physical Hamiltonian on \(\mathcal H_{\mathrm{phys}}\), put
\(E_{\mathrm{vac}}=\inf\operatorname{Spec}(A_{\mathrm{phys}})\), and assume
that \(E_{\mathrm{vac}}\) is an eigenvalue.  Let
\(P_0=\boldsymbol 1_{\{E_{\mathrm{vac}}\}}(A_{\mathrm{phys}})\), let
\(\mathfrak a_{\mathrm{phys}}\) be its closed quadratic form, and put
\(\mathcal H_{\mathrm{exc}}=(I-P_0)\mathcal H_{\mathrm{phys}}\).  This is a
reducing subspace.  Let \(\mathcal H_{\mathrm{aux}}\) be another Hilbert
space and let \(B\geq0\) be self-adjoint on it, with closed form
\(\mathfrak b\) satisfying
\begin{equation}
 \mathfrak b[v]\geq\gamma\|v\|_{\mathrm{aux}}^2
 \quad(v\in\operatorname{Dom}\mathfrak b)
 \label{hnsp-auxiliary-gap}
\end{equation}
for some \(\gamma>0\).

\begin{theorem}[Full-domain coercive transfer]
\label{hnsp-coercive-transfer}
 Suppose \(\mathcal H_{\mathrm{exc}}\ne\{0\}\) and a linear map
\begin{equation}
 R:\operatorname{Dom}(\mathfrak a_{\mathrm{phys}})
   \cap\mathcal H_{\mathrm{exc}}
 \longrightarrow\operatorname{Dom}(\mathfrak b)
 \label{hnsp-comparison-map}
\end{equation}
is defined on the entire physical excitation form domain and, for fixed
\(c_R,C_R>0\), satisfies
\begin{equation}
 \|Ru\|_{\mathrm{aux}}\geq c_R\|u\|_{\mathrm{phys}},
 \qquad
 \mathfrak b[Ru]\leq
 C_R\bigl(\mathfrak a_{\mathrm{phys}}[u]
             -E_{\mathrm{vac}}\|u\|^2\bigr).
 \label{hnsp-coercive-comparison}
\end{equation}
Then
\begin{equation}
 \inf\operatorname{Spec}
 \bigl(A_{\mathrm{phys}}|_{\mathcal H_{\mathrm{exc}}}\bigr)-E_{\mathrm{vac}}
 \geq\frac{\gamma c_R^2}{C_R}.
 \label{hnsp-transferred-gap}
\end{equation}
An intertwiner from only one gapped subspace cannot replace
\eqref{hnsp-comparison-map}.
\end{theorem}

\begin{proof}
For every vector in the domain of \eqref{hnsp-comparison-map},
\[
 C_R(\mathfrak a_{\mathrm{phys}}[u]-E_{\mathrm{vac}}\|u\|^2)
 \geq\mathfrak b[Ru]\geq\gamma\|Ru\|_{\mathrm{aux}}^2
 \geq\gamma c_R^2\|u\|_{\mathrm{phys}}^2.
\]
The variational characterization of the spectral bottom gives
\eqref{hnsp-transferred-gap}.  For the final assertion, take \(B=I\) on
\(\mathbb C\) and
\[
 A=0\oplus I\oplus M_x
 \quad\text{on}\quad
 \mathbb C\Omega\oplus\mathbb C\oplus L^2((0,1),dx),
\]
where \((M_x f)(x)=x f(x)\).  The isometry \(Jz=0\oplus z\oplus0\) obeys
\(AJ=JB\), yet \(A\) has unique vacuum \(\Omega\) and excitation spectral
bottom zero because of the multiplication summand.
\end{proof}

\subsection{Exact conclusion}

The determinant identity therefore supplies all of the following, exactly:
a canonical dimensionless affine Minkowski space, its Lorentz metric and
future cone, the spin and Poincar\'e actions, a natural time orientation, and
a dimensionless normal coordinate system.  After the independent choice of
one length/time dictionary this carrier becomes a spacetime model supporting a
standard continuum Klein--Gordon model, an exact finite periodic regulator,
and any chosen classical compact-group gauge theory.  The standard metrics
also select the compact \(K_P\cong SU(2)\) singlet--adjoint representation;
Theorem~\ref{hnsp-intrinsic-gauge-scalar-theorem} gives one exact gauging of
that representation.  Those are substantive physical constructions with
typed configuration spaces, states, actions, observables, spectra, and
limits.

The determinant normal-slice data alone do not select a preferred inertial
observer, an internal gauge interpretation, a coupling, a mass, a probability
measure, a quantum reconstruction, or a nonlinear spectral gap.  A separate
two-column Gram construction supplies a left \(SU(3)\)-action and the
positive momentum \(p_\alpha(Z)=\alpha Z^*Z\), but it retains the independent
parameters \(\alpha,c,m_0\) and does not dynamically select \(m_0>0\).
These limitations are exact rather than semantic: the inverse-dual candidate
cusp vector \eqref{hnsp-candidate-coordinates} is spacelike; no
positive quadratic energy is Lorentz invariant; the same host has the gapped
and gapless scalar realizations of Theorem~\ref{hnsp-free-field-gap}; the
massless gauge Gaussian gap closes as \(L\to\infty\); and entangled states
occur in both gapped and gapless Hamiltonians.  A claim about a physical mass
gap must therefore be proved from the physical spectrum or from a global
comparison controlling its entire excitation domain;
Theorem~\ref{hnsp-coercive-transfer} is one sufficient such criterion.
